
\documentclass{article} % For LaTeX2e
\makeatletter
\def\input@path{{iclr2027/}{./}}
\makeatother
\RequirePackage{iclr2027/fancyhdr,iclr2027/natbib}\usepackage{iclr2027_conference,times} % fancyhdr and natbib: the template's own copies (its \lhead is global, so the running header prints)

% Optional math commands from https://github.com/goodfeli/dlbook_notation.
\input{math_commands.tex}

\PassOptionsToPackage{hyphens}{url}\usepackage{hyperref}
\usepackage{url}
\usepackage{amsmath,amssymb}
\usepackage{booktabs}
\usepackage{multirow}
\usepackage{array}
\usepackage{xcolor}
\usepackage{graphicx}
\usepackage{enumitem}
\usepackage{tabularx}
\usepackage{longtable}
\usepackage{adjustbox}

% ---- the two mechanical markers ------------------------------------------------------
% \PLACEHOLDER{...}: a wave-15 / wave-16 result that is NOT written in. Its argument names
% the row and states that it may land as a negative. Wave 18's placeholder-fill greps it.
\newcommand{\PLACEHOLDER}[1]{{\color{red}\textbf{[PLACEHOLDER:}~#1\textbf{]}}}
% \TODO{...}: a gap this spoke could not fill from an artefact (listed in FINDINGS.md).
\newcommand{\TODO}[1]{{\color{orange}\textbf{[TODO:}~#1\textbf{]}}}
% scientific notation that the provenance checker can read back: \sci{7.522}{-16}
\newcommand{\sci}[2]{\ensuremath{#1\times10^{#2}}}
\newcommand{\Kstar}{\ensuremath{K^{\star}}}
\newcommand{\figslot}[3]{% #1 file stem in figures/, #2 caption, #3 label
  \begin{figure}[htb]\centering
  \IfFileExists{figures/#1.pdf}{\includegraphics[width=0.85\linewidth]{figures/#1.pdf}}%
  {\fbox{\parbox{0.9\linewidth}{\centering\small Figure pending (\texttt{figures/\detokenize{#1}.pdf} not yet
  available).}}}
  \caption{#2}\label{#3}\end{figure}}

\title{n-CLU: A Factorized Latent Dynamics Model for Deletable Compositional Generation}

% Authors must not appear in the submitted version. They are hidden as long as the
% \iclrfinalcopy macro remains commented out below.
\author{Anonymous authors}

%\iclrfinalcopy % Uncomment for camera-ready version, but NOT for submission.
\begin{document}

\maketitle

% ---------------------------------------------------------------------------------------
\begin{abstract}
A world model trusted with unseen tasks needs to go beyond generalizable predictions. Among other capabilities, dynamic tasks demand that (a) the model outputs are \emph{valid} under the task's judge (b) the model \emph{composes} new combinations of factors from the ones it has seen, and (c) the model can \emph{forget} specific factors when the environment changes permanently. Valid, compositional generation is the primary focus of this work with machine unlearning as a supplementary goal. We present the \emph{naive-Causal Learning Unit} (n-CLU),  and compare against naive baselines (an MLP, an LSTM, a transformer and a continuous-input MLP) on two simulated evaluation tasks: shape-rendering and maze-solving. The n-CLU is designed as a set of factorized, principal subspaces acting on the latent space of a frozen neural-network encoder, allowing composition of encoder-extracted features. The n-CLU is then read by a request-indexed address book, and finally rendered by a frozen decoder. Deleting a training row only downdates the statistic, so the trained n-CLU after deletion equals an n-CLU fitted without those rows. On a held-out set of the Shapes3D dataset,
the n-CLU composes new pairs and combinations of every factor. Against matched baselines trained on the same frozen encoder and decoder, the n-CLU produces $\sim 7$ more distinct, valid novel tuples per held-out combination than the strongest baseline, and beats every baseline on every evaluation set of an $n$-way combination task across multiple seeds. To test compositional generation in a harder, world-model aligned task, we created a maze-based benchmark (CMAZE-T) with composable factors and validity rules an agent must follow during navigation, based on the state of the local environment. We also release this small-scale benchmark for valid, compositional generation and plan to expand on it in future works. On CMAZE-T the n-CLU consistently beats or matches the strongest baseline across multiple seeds and we show both quantitative and qualitative results for the same. We finally motivate a Physics-designed extension to the n-CLU, and how it could provide clues towards designing better latent-dynamics based world models.
\end{abstract}

% =======================================================================================
\section{Introduction}
\label{sec:intro}
% =======================================================================================
Generating held-out combinations of factors is an active problem~\citep{du_compositional_2020,avidan_compositional_2022} in the AI community. Methods targeting composition, like the multilinear
class of algorithms (\citealp{tenenbaum_separating_2000},
TensorFaces~\citep{goos_multilinear_2002}, \citealp{vlasic_face_2005}, MLC-VAE~\citep{georgopoulos_multilinear_2020}) compose by factorizing a data tensor organized by every label combination and
rendering linearly. Other approaches based on the frozen-generator method (GANSpace~\citep{harkonen_ganspace_2020}, TCA~\citep{oldfield_tensor_2021}, PandA~\citep{oldfield_panda_2023} etc.) find editable directions in a pretrained generator without labels and without an address book.

On the other hand machine unlearning, that matches a retrain without the deleted rows, is also an active research area and a useful feature to have in a world model. $K$-means~\citep{ginart_making_2019} based methods, bilinear
matrix completion~\citep{xu_netflix_2023}, and ridge heads on frozen backbones~\citep{ceci_exact_2027} are common approaches, while principal-component eigen decomposition and generative models are relatively less explored. Recent generative-model unlearning attempts are approximate~\citep{kong_approximate_2022, kong_data_2023, moon_feature_2024, sun_generative_2025} applied to pretrained GANs and VAEs unlearn extracted features rather than input rows, and
are approximate by construction. SISA training~\citep{bourtoule_machine_2021} and compartmentalized diffusion models~\citep{golatkar_training_2024} on the other hand rely on retraining shards of a model and training multiple copies respectively and dropping the one that trained on the row to be deleted. To the best of our knowledge, the literature doesn't contain a unified solution for generation and input-row deletion.

Therefore, we present the naive-Causal Learning Unit (n-CLU), that combines multilinear factored generation, a request-indexed tensor head, and the downdate primitive, which follows from the summation form of~\citet{cao_towards_2015} into a single design acting as a Latent Dynamics Model (LDM). In this work we define an LDM as an algorithm that is either designed or learns to, enforce or extract the dynamical rules of a latent space parametrized by another independent learning algorithm. We later show how our n-CLU is a trivialization of a Physics-structured CLU, in-turn inspired by the CHLU presented in~\citet{jawahar2026chlu}. We believe LDMs like the CLU, as the name suggests, can provide actionable clues towards building AI models that enrich and leverage latent spaces better.

\subsection{World model properties}
\label{sec:why}

A world model~\citep{ha_world_2018,hafner_mastering_2025}, among other capabilities, must be able to generate valid, unseen states and transitions~\citep{valevski_diffusion_2025}.
An output is valid if the task's judge (regardless of whether a ground-truth answer exists) accepts it~\citep{vafa_evaluating_2024} and approximate similarity to a training item does not count.
Composition here, is defined as the model producing valid outputs for held-out combinations of factors it has seen only separately~\citep{zhao_toward_2022} or in a different composition. Deletion another useful feature and is defined as the algorithm matching a copy of itself that was never trained on an input or a factor, after that input or factor
is deleted. The three pull against each other: a lookup table can delete in this sense but cannot compose
anything not already stored; a learned generator composes but generally achieves deletion only by retraining
(the whole model, or a shard of it~\citep{bourtoule_machine_2021}); and
approximate unlearning~\citep{guo_certified_2023,golatkar_eternal_2020} is non-exact by construction. The n-CLU is a unified approach that combines some of these properties as explained below.

\subsection{Positioning and scope}
\label{sec:positioning}

The n-CLU is a primitive unit: one model with no harness built around it, in the sense in which MLPs, LSTMs~\citep{hochreiter_long_1997} or a transformers~\citep{vaswani_attention_2023} are AI primitives. We compare it only with other matched, naive units, an MLP, an LSTM, a transformer and a continuous-input MLP (Sec.~\ref{sec:object}). Systems built on top of units, such as sharded retraining, compositional diffusion and physics-informed variants of an MLP, are out of scope by definition. 


\subsection{Contributions}
\label{sec:contrib}

\begin{enumerate}
\item Compositional generation, against matched baselines on
  Shapes3D~\citep{3dshapes18,kim_disentangling_2019} and a world models-aligned maze dataset, CMAZE-T.
\item Deletion on the n-CLU and after
  a value, conjunction or item. Post-deletion, the n-CLU matches an n-CLU that was never trained on the deleted rows. However, we do not make a general claim about compositional refusal after input row deletion. 
\item A small-scale composition and deletion benchmark task (CMAZE-T) as a minor contribution. The benchmark consists of randomly generated maze grids, with colored doors, colored keys, tinted floors and a set of valid compositional rules that make maze solving less trivial and demand the algorithm to learn compositions to predict valid futures.
\end{enumerate}

% =======================================================================================
\section{Methods}
\label{sec:object}
% =======================================================================================

% PW2 item 3 (PLAN.md sec. 2): the Head's Methods opening, verbatim; LaTeX only (\citet = "Jawahar & Pierini (2026)").
\citet{jawahar2026chlu} introduced the CHLU as an LDM. We build on top of this idea and drop the physics-based design to construct our naive-CLU as a special case of the CHLU and test it as a factorized memory unit for compositional generation with additional properties such as deletion as a virtue of its design.

\subsection{Architecture and setup} 
The pre-requisites for the n-CLU as defined here are a pre-trained, frozen encoder that maps each training item to a latent space (this latent vector serves as the n-CLU input) and a pre-trained frozen decoder or downstream task head (takes n-CLU outputs and produces the output structure the task demands). The latent space is partitioned, without labels, into $K$ regions called wells by a balanced correlation partition. In each well the n-CLU keeps a fixed number of atoms, a hyperparameter: the mean of the well’s training latents and their leading principal directions. The atoms are not stored as such but derived from a retained per-well statistic, the count, sum and second-moment matrix of those latents, so that deleting rows amounts to subtracting them from the statistic and re-deriving the atoms. We define $n$ to be the number of factors composed for a given request at inference. Details of how $\Kstar$ is chosen are explained and tabulated in App.~\ref{app:kcurve}, Tab.~\ref{tab:kcurve3}).

\paragraph{The read step:} An inference request is a tuple of factor values. n-CLU contains a tensor-product address book,
defined over the training set, that maps the request to per-well coefficients. The read step renders
$x_k=U_k c_k$ in each well's orthonormal basis and the frozen decoder renders the image. The
read itself is request-only and no feature tensor enters it.
% SRC: w13-seeds-and-tables/TABLES.md
Labels therefore touch the model only at the address book. This is the interface split:
unsupervised context from a frozen encoder, a request-indexed head, and a frozen renderer, with no Cartesian-product
training requirement~\citep{wiedemer_compositional_2023,lachapelle_additive_2023,montero_lost_2022}.

\paragraph{The memory footprint:} Every n-CLU row shows four byte counts separately: (1) the atoms, (2) the covariance and count statistic, (3) the address book, and
(4) a coupling and hole table used for refusal. The 524,288-byte figure that recurs below is an arbitrary reference choice and the matched baselines carry about 131,072 (float32) learnable parameters. 

\paragraph{Factor relations under the hold-out draw:} The hold-out removes factor combinations from the training set at random. On some draws, for some factors, it removes every training
instance. A model that encodes, say color, by identity (like a lookup table) then has no training evidence for the
held-out instances of that color and cannot answer them, however it is trained. A model that is
also given the relation ``wall color equals tint color'' can: the relation is learned from the
other colors, where it is seen, and applies unchanged to the color whose rows were removed. This is
what the relational columns add (Fig.~\ref{fig:maze-wm}(c)), and it is why every seed is reported
separately for relational and non-relational models.
% ENDSRC

\paragraph{The judge:} On Shapes3D an external classifier, judges each rendered image; the metric is the number of \emph{distinct valid novel
factor tuples per held-out combination} (Tab.~\ref{tab:t7}). The results are reported over combinations with a margin of 1.0 tuple per held-out combination.

\paragraph{The retained statistic and deletion:} Along with the context the n-CLU also keeps per
well, the running covariance, mean and count of the fit rows. Deleting rows subtracts their
contribution from the statistic. The atoms are re-eigen decomposed from the subtracted (downdated)
statistic. The address book keeps its own retained statistic, the count and first
moment of its fit rows; deleting rows downdates it, and the book is refit from it.

% =======================================================================================
\section{Experiments}
\label{sec:experiments}
% =======================================================================================
The two datasets, Shapes3D and CMAZE-T are visualized Figs.\ref{fig:task-shapes} and \ref{fig:task-cmazet} respectively, and their setups are explained in the captions. 
\figslot{fig_task_shapes}{The Shapes3D task: (a) The six factors (floor hue, wall hue,
object hue, scale, shape, orientation), each row varying one factor of real dataset images with the
other five held at the request's values. (b) A request is a tuple of factor values; the answer is an
image rendered by the one frozen decoder, and the judge is the external classifier,
which must read back all six factors. (c) The two-way hold-out on
object hue $\times$ shape (d) A six-factor combination hold-out example.}{fig:task-shapes}

\figslot{fig_task_cmazet}{The CMAZE-T task: (a) One generated maze, drawn by the engine:
walls, floor tinted in one of four colors, colored keys (diamonds),
closed and opened doors, the goal (a cross), and the agent: a black disc whose center shows the key in
its pocket, white when the pocket is empty (it holds at most one key).
Rule A: a closed door of color $c$ needs key $c$ to open and consumes the key after. Rule B: stepping on floor tinted $c$ is forbidden
while holding key $c$. Rule A is applied before rule B. (b)--(d) One transition each: the inference request is
five words (\texttt{w\_agent}, \texttt{w\_key}, \texttt{w\_door}, \texttt{w\_tint}, \texttt{w\_dir})
and the answer is four facts (\texttt{legal}, \texttt{disp}, \texttt{event}, \texttt{pocket}),
computed by the engine and checked against an independent stepper. (b) A legal move onto floor of
another tint. (c) An illegal move under rule B: the agent does not move, so after equals before and
the verdict is carried by the mark and by \texttt{legal}~$=0$. (d) The same pocket and tint as (c) with a closed door of the
pocket's color: the door consumes the key first, so the move is legal. Panels (c) and (d) differ
only in the door, which is the pocket $\times$ door $\times$ tint interaction held out in this example.}{fig:task-cmazet}
% ENDSRC

\paragraph{Baselines:}

Learned baselines (MLP, LSTM~\citep{hochreiter_long_1997}, transformer~\citep{vaswani_attention_2023}, and a
continuous-input MLP) are matched to the n-CLU's parameter count, and fitted from the same frozen features. At $\Kstar$ the fitted baselines carry the parameter counts of Tab.~\ref{tab:t7}.
A lookup table
of written items is a non-learned comparison floor in our case. The four controls (landscape controls) that help isolate the contribution of the n-CLU over the frozen encoder-decoder are: \emph{landscape removed} (atoms zeroed), \emph{landscape shuffled} (atom values permuted within each well), \emph{shuffled across wells}
(well blocks dealt to other wells) and \emph{random atoms} (random atoms in place of the fitted ones).

\paragraph{$n$-way composition:}

This experiment checks the n-CLU's composition performance as more factors are requested jointly in a single inference pass out of different hold outs. Tab.~\ref{tab:ladder} shows the n-CLU at $\Kstar$ on every $n$ and
seed beside matched baseline sets on the same combinations.
The landscape controls at $\Kstar$ collapse on every seed as seen in Tab.~\ref{tab:t7}, showing the n-CLU is not being by-passed by the encoder-decoder. 

\paragraph{Deletion:}

Three deletions are run at $\Kstar$, on its fit pool (Tab.~\ref{tab:rivals-del-seeds}), at three % SRC: w20-shapes-data-matched/FINDINGS.md
% PW2-5 moved -> Tab. del caption (fit pool 131,072 rows): 131,072; SRC: w20-shapes-data-matched/FINDINGS.md
seeds. \emph{Value}: every fit row carrying red. \emph{Conjunction}: every fit row carrying red
\emph{and} sphere. \emph{Item}: a random 1\% of fit rows, which is
also the control for refit instability. \emph{Forgetting $c$} is three operations:
downdate the n-CLU's context; downdate the address book's statistic and refit the book from it; and punch the
hole for $c$ (Sec.~\ref{sec:hole}) for post-deletion refusal.

% =======================================================================================
\section{Results}
\label{sec:results}
% =======================================================================================

\subsection{Composition on Shapes3D}
\label{sec:shapes-t7}

Tab.~\ref{tab:t7} shows the full composition results on Shapes3D. Row 3 holds out 60 pair combinations and row 4 holds out
49 six-factor combinations; the n-CLU is compared with each baseline on all 109 combinations.

\begin{table}[htb]
\centering\small
\caption{Composition at $\Kstar=32$ against matched baselines on Shapes3D, $3$
seeds (0, 1, 2). The n-CLU % SRC: w13-seeds-and-tables/TABLES.md
and every baseline (no. of training steps in brackets) is matched to $\sim130k$ parameters.  ``Paired $\Delta$'' is the n-CLU minus
that baseline over all 109 combinations. LACE~\citep{nie_controllable_2021} and PoE~\citep{hinton_training_2002} are in % SRC: w12b-nr6-and-kstar/NR6-T7.md
Tab.~\ref{tab:gentable-kstar}. The n-CLU's row-3 and row-4 means are 32.350~$\pm$~0.815 and 67.177~$\pm$~1.320 % SRC: w12b-nr6-and-kstar/NR6-T7.md
($\pm$ the sample standard deviation over the three seeds).}
\label{tab:t7}
\setlength{\tabcolsep}{3pt}
\begin{adjustbox}{max width=\linewidth}\begin{tabular}{lcccccc}
\toprule
arm (optimiser steps) & seeds & all 109 combinations & (60 pairs held out) & (49 $n=6$ held out) & paired $\Delta$, n-CLU $-$ arm & combinations won / tied / lost \\ % SRC: w12b-nr6-and-kstar/NR6-T7.md
 & & seed 0 / 1 / 2 & 3-seed mean & 3-seed mean & seed 0 / 1 / 2 & seed 0 $\cdot$ 1 $\cdot$ 2 \\
\midrule
% SRCBLOCK: w20-shapes-data-matched/SHAPES-MATCHED.md; w20-shapes-data-matched/SHAPES-MATCHED.json; w20-shapes-data-matched/rivals/test30x/seed0/mlp/declared.json
\textbf{n-CLU, $\Kstar=32$ (closed form)} & 3 & \textbf{47.495 / 47.404 / 49.119} & \textbf{32.350} & \textbf{67.177} & --- & --- \\
MLP (20,480) & 3 & 36.706 / 37.569 / 43.174 & 25.394 & 55.993 & & \\
MLP (204,800)$^\ast$ & 3 & 40.009 / 38.807 / 44.101 & 26.417 & 58.796 & +7.486 / +8.596 / +5.018 & 76/11/22 $\cdot$ 84/9/16 $\cdot$ 69/27/13 \\
MLP (614,400; diagnostic) & 3 & 41.477 / 36.541 / 43.495 & 25.006 & 59.483 & +6.018 (seed 0) & 67/14/28 (seed 0) \\
LSTM (20,480)$^\ast$ & 3 & 33.972 / 32.358 / 40.697 & 23.467 & 50.626 & +13.523 / +15.046 / +8.422 & \\
LSTM (204,800) & 3 & 33.688 / 32.541 / 39.853 & 22.994 & 50.503 & & \\
transformer (20,480)$^\ast$ & 3 & 25.560 / 27.789 / 38.651 & 20.050 & 43.667 & +21.936 / +19.615 / +10.468 & \\
transformer (204,800) & 3 & 25.303 / 27.239 / 37.844 & 18.783 & 44.020 & & \\
continuous MLP (20,480) & 3 & 17.917 / 16.541 / 26.844 & 11.722 & 31.102 & & \\
continuous MLP (204,800)$^\ast$ & 3 & 28.532 / 32.936 / 42.596 & 22.589 & 49.503 & +18.963 / +14.468 / +6.523 & \\
continuous MLP (614,400; diagnostic) & 3 & 27.284 / 34.009 / 38.706 & 21.894 & 47.340 & & \\
lookup table (floor) & 3 & 8.229 / 5.073 / 9.752 & 5.167 & 10.769 & & \\
\midrule
four landscape controls & 3 & 0.000 (all) & 0.000 (all) & 0.000 (all) & & \\
% ENDSRC
\bottomrule
\end{tabular}\end{adjustbox}
\end{table}

\paragraph{Qualitative results:} Figs.~\ref{fig:shapes-per-model} and~\ref{fig:shapes-baseline-wins}
show single held-out requests from Tab.~\ref{tab:t7} as images, at seed 0 and $\Kstar=32$, against four baselines the lookup table. Every model's output on an inference request is drawn through the same frozen decoder, and each panel carries the judge's verdict on that image. The CLU gets most combinations right when atleast one baseline fails. When the n-CLU is incorrect, its image is often wrong in one factor only (most often the scale or the orientation, less often the shape), which is potentially a systematic failure mode at the current configuration. Requests on which a baseline succeeds and the n-CLU fails exist for every baseline and are not confined to the capsule 
(Fig.~\ref{fig:shapes-baseline-wins}).

\subsection{The $n$-way composition}
\label{sec:ladder}

\begin{table}[t]
\centering\small
\caption{$n$-way composition on Shapes3D, $3$ seeds. Distinct valid novel tuples per combination reported, per seed (0 / 1 / 2). n-CLU outperforms all baselines at all combination-seed pairs.}
\label{tab:ladder}
\setlength{\tabcolsep}{2.6pt}
\begin{adjustbox}{max width=\linewidth}\begin{tabular}{lcccccc}
\toprule
arm & seeds & $n=2$ & $n=3$ & $n=4$ & $n=5$ & $n=6$ \\
\midrule
% SRCBLOCK: w12b-nr6-and-kstar/KSTAR-LADDER.md
\textbf{n-CLU, $\Kstar$} & 3 & \textbf{31.800 / 32.500 / 33.633} & \textbf{60.150 / 62.783 / 61.908} & \textbf{65.242 / 65.817 / 67.200} & \textbf{65.642 / 66.817 / 67.283} & \textbf{65.908 / 66.250 / 67.650} \\
\quad 3-seed mean & 3 & 32.644 & 61.614 & 66.086 & 66.581 & 66.603 \\
% ENDSRC
% n-CLU, $n=6$ on $n=3$'s fit pool & 3 & \multicolumn{4}{l}{TEST cell \emph{derived, not measured}: the $n=3$ n-CLU cell, by construction (text)} & \emph{60.150 / 62.783 / 61.908, derived} \\ % SRC: w15-coverage-control/FINDINGS.md
\midrule
\multicolumn{7}{l}{\emph{data-matched set: each rung's own 131,072 rows; 20,480 steps, and 204,800 for the MLP and the continuous MLP}} \\ % SRC: w20-shapes-data-matched/FINDINGS.md; w21-ladder-rivals-32k/CAPTION.md
% SRCBLOCK: w20-shapes-data-matched/SHAPES-MATCHED.md; w20-shapes-data-matched/SHAPES-MATCHED.json
MLP (20,480) & 3 & 22.450 / 25.200 / 29.267 & 43.558 / 46.742 / 51.592 & 45.408 / 48.642 / 57.300 & 51.192 / 52.825 / 58.842 & 49.525 / 54.692 / 58.867 \\
MLP (204,800) & 3 & 27.167 / 25.583 / 29.783 & 47.258 / 46.833 / 57.200 & 50.517 / 50.017 / 60.075 & 53.642 / 53.275 / 60.767 & 54.183 / 52.200 / 60.033 \\
LSTM (20,480) & 3 & 20.550 / 22.383 / 27.150 & 36.400 / 37.233 / 40.908 & 42.492 / 39.958 / 53.517 & 46.492 / 47.442 / 55.692 & 47.758 / 46.275 / 55.417 \\
transformer (20,480) & 3 & 16.417 / 17.183 / 25.750 & 26.083 / 25.192 / 30.550 & 31.792 / 33.717 / 48.800 & 37.142 / 37.458 / 52.475 & 36.592 / 37.300 / 52.942 \\
continuous MLP (20,480) & 3 & 10.383 / 11.517 / 11.383 & 16.008 / 21.225 / 29.167 & 21.000 / 24.642 / 36.183 & 21.817 / 32.583 / 37.550 & 19.817 / 29.517 / 39.875 \\
continuous MLP (204,800) & 3 & 21.333 / 21.950 / 24.867 & 24.125 / 41.300 / 39.808 & 29.883 / 45.158 / 44.883 & 40.192 / 49.858 / 54.258 & 48.633 / 49.400 / 49.525 \\
\textbf{n-CLU $-$ strongest, paired} & 3 & +4.633 / +6.917 / +3.850 & +12.892 / +15.950 / +4.708 & +14.725 / +15.800 / +7.125 & +12.000 / +13.542 / +6.517 & +11.725 / +11.558 / +7.617 \\
\quad 3-seed mean & 3 & +5.133 & +11.183 & +12.550 & +10.686 & +10.300 \\
% ENDSRC
\midrule
\multicolumn{7}{l}{\emph{n-CLU-context set: the rung's first 32,768 rows, exactly the n-CLU's context rows; 5,120 steps, ten times 51,200}} \\ % SRC: w21-ladder-rivals-32k/CAPTION.md
% SRCBLOCK: w21-ladder-rivals-32k/TABLES.md
MLP (5,120) & 3 & 22.07 / 25.13 / 25.90 & 35.79 / 41.36 / 47.81 & 39.32 / 43.39 / 50.93 & 43.64 / 47.91 / 53.64 & 45.77 / 48.77 / 54.26 \\
MLP (51,200) & 3 & 25.15 / 24.82 / 29.42 & 47.14 / 47.42 / 55.84 & 49.10 / 49.09 / 58.07 & 54.52 / 56.22 / 60.59 & 54.61 / 53.76 / 60.36 \\
LSTM (5,120) & 3 & 22.42 / 20.67 / 26.55 & 30.52 / 35.98 / 43.42 & 41.54 / 39.98 / 50.76 & 46.10 / 44.67 / 52.17 & 45.64 / 43.67 / 55.12 \\
LSTM (51,200) & 3 & 22.58 / 22.42 / 27.05 & 36.39 / 37.67 / 45.37 & 42.68 / 41.33 / 53.71 & 45.95 / 45.62 / 55.12 & 48.17 / 46.34 / 55.36 \\
transformer (5,120) & 3 & 13.87 / 17.50 / 25.17 & 27.98 / 24.17 / 29.76 & 32.04 / 31.27 / 47.18 & 35.08 / 39.42 / 53.47 & 36.87 / 37.83 / 52.74 \\
transformer (51,200) & 3 & 16.85 / 16.58 / 23.87 & 26.02 / 25.63 / 35.48 & 30.60 / 32.56 / 47.44 & 36.40 / 36.44 / 52.02 & 35.44 / 36.19 / 52.42 \\
continuous MLP (5,120) & 3 & 3.95 / 6.23 / 10.03 & 6.88 / 11.33 / 14.18 & 6.10 / 8.65 / 18.28 & 9.78 / 17.45 / 23.59 & 9.15 / 13.01 / 23.92 \\
continuous MLP (51,200) & 3 & 16.88 / 17.38 / 22.32 & 24.69 / 31.81 / 35.46 & 20.59 / 27.65 / 46.14 & 38.71 / 32.98 / 45.76 & 33.28 / 34.03 / 45.79 \\
\textbf{n-CLU $-$ strongest, paired} & 3 & +6.65 / +7.37 / +4.22 & +13.01 / +15.37 / +6.07 & +16.14 / +16.73 / +9.13 & +11.12 / +10.60 / +6.69 & +11.30 / +12.49 / +7.29 \\
\quad 3-seed mean & 3 & +6.08 & +11.48 & +14.00 & +9.47 & +10.36 \\
% ENDSRC
\bottomrule
\end{tabular}\end{adjustbox}
\end{table}

n-CLU outperforms all baselines at
every $n$ as seen in Tab.~\ref{tab:ladder}. However, the rise from 2-way to 6-way is not to be inferred, because rungs 3--6 probe overlapping combinations and differ strictly only in fit-pool coverage. What varies is how
coarsely the hold-out set is carved, and with it the held-out row count. Regardless, on every rung and every seed the n-CLU's score in Tab.~\ref{tab:ladder} exceeds every baseline's
score (Fig.~\ref{fig:ladder}).

\begin{figure}[htb]\centering
\includegraphics[width=0.9\linewidth]{figures/rendered/rc_pw1_shapes_per_model.pdf}
\caption{Each model shown succeeding and failing on held-out combinations of Shapes3D, seed 0. The model names on the left span 2 rows each, one that model got right and one wrong. Of the n-CLU's 511 invalid renders, 497 are wrong in one factor only: scale 329, orientation 125, shape 43.}
\label{fig:shapes-per-model}
\end{figure}

\subsection{Deletion on Shapes3D}
\label{sec:shapes-del}

% SRCBLOCK: w20-shapes-data-matched/SHAPES-MATCHED.md
\noindent After deletion the n-CLU matches an n-CLU that was never fitted on the deleted rows, and its address
book matches a book refit without them: over every scope and seed the context's largest mode
deviation and the relative deviation of the book's rendered features are at floating-point roundoffs, and the book's statistic does not deviate.

\paragraph{Conjunction deletion and the shared value:} This is a
measured limitation of the n-CLU. Deleting the red sphere lowers composition on the surviving combinations
by about six tuples per held-out combination, on all three seeds (Fig.~\ref{fig:deletion}). At seed 0 the loss is concentrated where the deleted rows tied red to a shape: the combinations that request red with another shape fall by about thirty each, those that request a sphere in another color by about one, and the others don't drop. The same pipeline deleting a random 1\% of items only drops on remaining requests by at most about a tenth of a tuple, so the loss is the deleted context. The drop is what a book
fitted without the red sphere's rows no longer knows about red. The targeting of a deletion is a function of rank/width, and it disappears where the memory footprint
is smallest. This is why refusal needs its own mechanism (Sec.~\ref{sec:hole}). Further results and discussions on deletion for the baselines and comparing that with the n-CLU are presented in App.~\ref{app:delbase}).


\paragraph{Encoder adaptation and generative baselines}
\label{sec:adapt}\label{sec:bard}

The frozen encoder is the primary object. Light adaptation at $\Kstar$ with equal steps and
images grows the n-CLU's margin over the best learned baseline on all three seeds under an
anchored full unfreeze and on two
of three under a low-rank adapter (Fig.~\ref{fig:adapt}). % SRC: w13-

\subsection{CMAZE-T results and relational coding}
\label{sec:maze-wm}

To isolate the compositional generation ability of the models, the benchmark is split into single step action prediction for a single agent given a complex rule-based condition on the answer.

\figslot{fig_maze_conditions}{CMAZE-T: the n-CLU against the best matched baseline under
four input choices, and six seeds without and with relational coding.(a)~Unseen $n$-way validity
under the engine's judge, one row per input choice: the full word design, the
observation, and words plus observation for the n-CLU and for the n-CLU-REL. Every seed of a row is drawn: circles are the n-CLU (n-CLU-REL in the last row),
diamonds the row's baseline, and a vertical tick marks the mean over the row's seeds. The baseline of
the first three rows is the matched Transformer with the highest mean over the row's seeds, taken over
every family and both step limits. The baseline of the n-CLU-REL
row is Transformer-REL, the Transformer given the same relational columns (hollow diamonds). (b)~Words plus observation, test, seeds 0--5, without relational coding: per
seed, the n-CLU, the Transformer at $10\times$ training steps of the third row, and the highest of the
four controls on that seed. (c)~The same as (b) but with relational coding: the n-CLU-REL, the Transformer-REL of
the last row, and the highest control.}{fig:maze-wm}

% SRCBLOCK: w21-maze-t-seed4/FINDINGS.md; w21-maze-t-seed4/TABLES.md; w21-words-full-test/FINDINGS.md; w19-maze-t-symmetric/FINDINGS.md
\paragraph{CMAZE-T with symmetric inputs:}
On the full word design given % pw1-draft-sweeps (item 2): the base design / base n-CLU was Config A1 and n-CLU-REL was Config CLU-REL-A10 in the w16/w19/w21 records
to every arm it roughly matches the best baseline (transformer) on all three TEST seeds, as seen in Fig.~\ref{fig:maze-wm}(a). On words
plus observation, on six seeds, against the
transformer it has three wins, two ties and
one loss. The 10-atom relational design (n-CLU-REL, below) scores every row of both tuples valid. 
All controls stays below the n-CLU on every seed.

\paragraph{Relational request coding:} 
The relational coding is an additional set of inputs that show a map between related words, such as "tinted floor" and "blue key" are both related since they contain color as a main factor while "move East" has no color association. Fig.\ref{fig:maze-wm}(c) shows that adding the relational coding improves valid actions more than not having them but, they equally lift the transformer arm as well. Regardless, the n-CLU and n-CLU-REL are ahead of the best baseline marginally across $6$ seeds.

\figslot{fig_cmazet_families}{Held-out CMAZE-T transitions, one per rule family, with every
model's next state (words plus observation, test).}{fig:cmazet-families}


\paragraph{The CMAZE-T panel figures:} Figs.~\ref{fig:cmazet-families}, \ref{fig:cmazet-per-arm}--\ref{fig:cmazet-door-tint}
(App.~\ref{app:mazewm}) draw single held-out transitions with every model's next state. Each panel
is a $5\times5$ crop of the maze centered on the agent's cell at $s_t$. Walls are dark, tinted cells
pale in their color, doors framed in their color with a bar across while closed, untaken keys are
diamonds, the goal a cross, and the agent a black disc whose center shows the color of the key in its
pocket (white when the pocket is empty). The first column is $s_t$ with the action (an arrow, or a
ring for pick-up), the second the engine's $s_{t+1}$, and each further column the $s_{t+1}$ that a
model's emitted outcome produces. The engine's judge does not score the picture: it scores four
emitted facts, legality, the move, the key event and the pocket after the move. A picture need
not show all four. So a strip under every $s_{t+1}$ prints the emitted facts (under the engine's
column, the true ones). Its first line reads legal or illegal. Its second line shows the move
($\uparrow$ $\rightarrow$ $\downarrow$ $\leftarrow$, or a dot for staying), the key event ($+$ and a
color for a key picked up, $-$ and a color for a key used at a door, blank for none) and the pocket
after the move (the disc; white is empty). An illegal emission leaves the agent in place and is scored
on its pocket alone, so only the pocket is shown. The tick or cross on a panel is the judge's verdict
on all four facts, so two panels that show the same picture and the same strip never carry opposite
marks.

% =======================================================================================
\section{Discussion}
\label{sec:limits}
% =======================================================================================

In this work we present n-CLU and demonstrate its compositional generation and deletion, matching an n-CLU never fitted on the deleted rows, at a stated partition.

The n-CLU we evaluate is the zero-dynamics limit of the CHLU~\citep{jawahar2026chlu}: the same context, wells and read, with the settling loop's step count set to zero. Inside a well this is an
algebraic identity: the n-CLU's energy is the distance to the atom span, the request's basin sits
at an address inside that span, so both minima coincide and an in-well settle has nowhere to go
(derivation in App.~\ref{app:clue}).

We show performance on
Shapes3D that scores above all matched baselines by $\sim$seven new tuples per held-out combination over multiple seeds. 

We also create and provide a small scale dataset CMAZE-T targeted at world models and their valid, composition abilities. CMAZE-T is the first instrument of a planned CMAZE-suite on the same
engine and judge. The transition task will be extended to observation generation, to route planning and to the further tasks already prototyped, so that CMAZE as a whole becomes a useful benchmark for
the world-models community. On CMAZE-T n-CLU's lead over matched
baselines is small and seed-dependent. 

The main limitations we plan to expand upon are improving the memory footprint required for deletion, expanding the CMAZE-suite to more complex, long-horizon cases, and evaluating the n-CLU on real world-model relevant data. 

The follow-up is a full evaluation of the settling (physics) CLU on the three properties of
Sec.~\ref{sec:why}, same judges, memory footprint and matching, with this paper's n-CLU as its zero-dynamics
ablation. Our future work also includes improving the memory footprint when moving to the physics version. We also intend to use the n-CLU as a diagnostic tool to understand how best to construct the latent physics-CLU.


% ---------------------------------------------------------------------------------------
\subsubsection*{AI use statement}
In this work agentic coding tools were used to draft and automate the code based experiments. The code was verified using previously human-authored code harnesses used in prior research, by 2 of the authors who worked on the same datasets presented here (as a unit-test style quality control). The results generated were evaluated and validated by human authors. Generative AI was used to draft the initial skeleton of this manuscript with 2 human authors responsible for line-by-line major revisions and edits of the text and plots and 2 other authors provided editorial passes on the human-revised version, blind to the initial AI generated draft. We take responsibility for the final content of this work, including text, claims or artifacts produced with the aid of generative AI.

Below, we share a list of common research subtasks along with the necessity of disclosing them in the paper manuscript:

Generate synthetic data sets, help develop theoretical models or conceptual frameworks, formulate mathematical claims, provide critical ingredients for proving mathematical claims, assist in the writing of proofs, implement methods, clean and reformat dataset, support qualitative and thematic data analysis, create or modify scientific figures or images, create or edit software code, creation of artifacts, draft parts of a research paper, summarize or analyze existing literature, brainstorming, sourcing/searching for information, identify relevant literature, propose keywords for a research paper.

\subsubsection*{Ethics statement}
The study uses a public synthetic image dataset (Shapes3D) and procedurally generated mazes; it
involves no human subjects and no personal data. Deletion is relevant to data-protection
requests; we claim only that the n-CLU and its address book match those never fitted on the deleted
rows (the frozen encoder and decoder are outside the object), no certified unlearning is to be inferred, and
we state in Sec.~\ref{sec:shapes-del} a case in which deleting a conjunction changes what the
model answers about the values that survive it.

\subsubsection*{Reproducibility statement}
Each number in this paper is read from the recorded result of a declared run, and nearly every
one carries its provenance beside it in the source. The verbatim re-check against the cited records
is partial: an automated check covers every number of three groups of runs and, overall, 1,308 of % SRC: ../../outputs/w23-fold/FINDINGS.md
the 2,293 sourced number tokens of an earlier version of this draft, and it checks that a number is % SRC: ../../outputs/w23-fold/FINDINGS.md
present in its record, not what the number means. The object, judge, estimand and memory footprint are
defined in Sec.~\ref{sec:object}; every table states its seeds, condition and pins. The anonymised code archive:\url{https://anonymous.4open.science/r/n-clu-F746/}.

% ---- references (cite-tagger owns references.bib; this draft carries no \cite yet) --------
\IfFileExists{refs_clu_iclr.bib}{%
\bibliography{refs_clu_iclr}
\bibliographystyle{iclr2027/iclr2027_conference}}{} % (no result number on this line: wave labels / file names only)

% =======================================================================================
\clearpage
\appendix
\section*{Appendices}

% ---------------------------------------------------------------------------------------


% ---------------------------------------------------------------------------------------

\section{The $K$-curve: composition and the memory footprint against $K$}
\label{app:kcurve}


\figslot{fig_crossover}{The $K$-sweep validation, 3 seeds at every $K$. (a)~Validation
composition against $K$: the n-CLU's mean over seeds 0, 1 and 2 with a band of one seed standard
deviation, each seed marked, and the highest of the four controls (landscape removed, landscape
shuffled, shuffled across wells, random atoms). (b)~The memory footprint in its four parts (atoms,
covariance and count, address book, coupling and hole table; symmetric-triangle convention). The values are printed in Tab.~\ref{tab:kcurve3}.}{fig:crossover}

\begin{table}[h]
\centering\scriptsize
\caption{The generation table at $\Kstar=32$. Outperforms LACE~\citep{nie_controllable_2021}, PoE~\citep{hinton_training_2002}). Caveats mentioned below.}
\label{tab:gentable-kstar}
\setlength{\tabcolsep}{2.4pt}
\begin{adjustbox}{max width=\linewidth}\begin{tabular}{lcccp{4.2cm}}
\toprule
arm & seeds & aggregate, seed 0 / 1 / 2 & mean & deployed bytes (convention) \\
\midrule
% SRCBLOCK: w15-lace-poe-rows/FINDINGS.md
\textbf{$\Kstar=32$} & 3 & \textbf{0.5718 / 0.9977 / 0.6806} & \textbf{0.7500 $\pm$ 0.1807} & 2,145,738 / 2,153,018 / 2,155,618 B, four-part, never netted (4.09--4.11$\times$ ref) \\
\quad landscape removed & 3 & 0.0556 / 0.0556 / 0.2778 & 0.1296 & same object, atoms zeroed \\
\quad landscape shuffled & 3 & 0.1481 / 0.1389 / 0.0602 & 0.1157 & same bytes \\
\quad shuffled across wells & 3 & 0.1065 / 0.0764 / 0.1435 & 0.1088 & same bytes \\
\quad random atoms & 3 & 0.1458 / 0.2523 / 0.1597 & 0.1860 & same bytes \\
earlier n-CLU's read & 3 & 0.5903 / 0.6065 / 0.5880 & 0.5949 $\pm$ 0.0082 & 524,288 B (n-CLU modes) \\
n-CLU removed & 3 & 0.1111 / 0.1389 / 0.1111 & 0.1204 & 524,288 B \\
nearest-prototype ceiling & 3 & 0.8588 / 0.8542 / 0.8819 & 0.8650 $\pm$ 0.0122 & 475,136 B (selected under a retired byte limit, declared) \\
LACE (\texttt{ode}) & 3 & 0.2940 / 0.3241 / 0.2917 & 0.3032 $\pm$ 0.0148 & not byte-matched \\
PoE (\texttt{ebm\_compose}) & 3 & 0.1759 / 0.1875 / 0.2269 & 0.1968 $\pm$ 0.0218 & 14,226,660 B (27.1$\times$ ref) \\
% \midrule
% $\Kstar$ $-$ earlier n-CLU's read (S2) & 3 & $-0.0185$ / $+0.3912$ / $+0.0926$ & $+0.1551$ & \textsc{fail} (seed 0) \\
% $\Kstar$ $-$ nearest-prototype ceiling (S3) & 3 & $-0.2870$ / $+0.1435$ / $-0.2014$ & $-0.1150$ & \textsc{lose} \\
% ENDSRC
\bottomrule
\end{tabular}\end{adjustbox}
\end{table}


\paragraph{LACE and PoE:}
LACE~\citep{nie_controllable_2021} and PoE~\citep{hinton_training_2002} were scored beside the
learned baselines on rows 3--4 and on every ladder rung (Tab.~\ref{tab:gentable-kstar}), 3 seeds. No
learned-baseline column was replaced. They do not beat the
n-CLU on rows 3--4, nor anywhere on the ladder: on rows 3--4 the $\Kstar$ n-CLU scores 47.50 / % SRC: w15-lace-poe-rows/FINDINGS.md
47.40 / 49.12 novel valid tuples per held-out combination against LACE 0.165 / 0.073 / 0.229 and PoE 0.009 / % SRC: w15-lace-poe-rows/FINDINGS.md
0.000 / 0.000, and the n-CLU wins 109/109 combinations against each, on every seed; on every rung it % SRC: w15-lace-poe-rows/FINDINGS.md
wins every combination against both on every seed (60 of 60 at $n=2$, 120 of 120 at $n\ge3$). % SRC: w15-lace-poe-rows/FINDINGS.md
Their settings are their own and are not byte-matched: LACE 4,073,663 classifier parameters % SRC: w15-lace-poe-rows/FINDINGS.md
(31.08$\times$ the 524,288 B reference; 63.10$\times$ counting its whitening and mean), PoE six % SRC: w15-lace-poe-rows/FINDINGS.md
experts with 3,556,665 parameters (27.135$\times$), each fitted on 32,256 rows with 512 validation % SRC: w15-lace-poe-rows/FINDINGS.md; w15-lace-poe-rows/GENTABLE-KSTAR.json
rows from the rung's own pool and 36 selection trials. % SRC: w15-lace-poe-rows/FINDINGS.md
\emph{This is LACE and PoE as instantiated in this feature space at their own declared settings, it is not a refutation of LACE given there is no generative baseline.}


\begin{table}[h]
\centering\small
\caption{The $K$-curve at $3$ seeds on every row, on the validation set.}
\label{tab:kcurve3}
\setlength{\tabcolsep}{2.4pt}
\begin{adjustbox}{max width=\linewidth}\begin{tabular}{rcccccccc}
\toprule
$K$ & seeds & rank/width & val comp. mean & seed sd & drop vs $K{=}5$ & four-part B mean & seed sd B & strict check \\
\midrule
% SRCBLOCK: w15-crossover-seeds/CROSSOVER-3SEED.md
5 & 3 & 0.1561--0.1565 & 0.8953 & 0.0065 & +0.0000 & 4,118,034 & 0 & pass \\
8 & 3 & 0.2500--0.2500 & 0.9037 & 0.0197 & $-0.0085$ & 2,959,602 & 0 & pass \\
10 & 3 & 0.3122--0.3137 & 0.9014 & 0.0211 & $-0.0062$ & 2,606,754 & 0 & pass \\
12 & 3 & 0.3743--0.3765 & 0.9082 & 0.0141 & $-0.0129$ & 2,393,714 & 0 & pass \\
16 & 3 & 0.5000--0.5000 & 0.9111 & 0.0203 & $-0.0158$ & 2,177,330 & 0 & pass \\
20 & 3 & 0.5294--0.6275 & 0.9178 & 0.0285 & $-0.0226$ & 2,098,879 & 2,864 & \textbf{fail} \\
24 & 3 & 0.4353--0.7529 & 0.9223 & 0.0235 & $-0.0270$ & 2,086,677 & 7,060 & pass \\
28 & 3 & 0.3425--0.8767 & 0.9267 & 0.0276 & $-0.0314$ & 2,117,183 & 9,885 & pass \\
32 & 3 & 0.2656--0.9844 & 0.9290 & 0.0259 & $-0.0338$ & 2,140,538 & 3,640 & pass \\
% ENDSRC
\bottomrule
\end{tabular}\end{adjustbox}
\end{table} % SRC: w13-seeds-and-tables/TABLES.md


\figslot{fig_ladder}{The ladder at $\Kstar$ against matched baselines on Shapes3D, across 3 seeds showing the distinct
valid new tuples created per held-out combination against $n$, the number of factors jointly held out.}{fig:ladder}


\begin{figure}[t]\centering
\includegraphics[width=\linewidth]{figures/rendered/rc_pw1_shapes_baseline_wins.pdf}
\caption{Held-out requests on which a baseline succeeds and the n-CLU fails (seed 0).}
\label{fig:shapes-baseline-wins}
\end{figure}


% ---------------------------------------------------------------------------------------
\section{Deletion: Extended Discussion}
\label{app:delbase}

% ---------------------------------------------------------------------------------------

\figslot{fig_deletion}{Deletion at $\Kstar=32$, rows 3--4, test set, fit pool
131,072 rows, seeds 0, 1 and 2. Deleting a value (red), a conjunction (red sphere) or a random 1\% % SRC: w20-shapes-data-matched/FINDINGS.md
of items downdates the n-CLU's context, downdates the address book's statistic and refits the book
from it, and punches the hole. (a)~Composition on the remaining held-out
combinations (89 for the value, 109 for the others), after minus before, per seed, with its 95\% % SRC: w20-shapes-data-matched/SHAPES-MATCHED.json
bootstrap interval.
(b)~The n-CLU after deletion minus each baseline retrained without the same rows of the same pool,
with 95\% intervals: above zero for every baseline and seed on the value and item deletions; on red % SRC: w20-shapes-data-matched/SHAPES-MATCHED.json
sphere the interval includes zero against the MLP at seeds 1 and 2 and against the LSTM at seed 2.}{fig:deletion}


\subsection{Deletion against retrained baselines}
\label{sec:del-comparators}

A learned baseline's only route to a model never fitted on the rows is to retrain from scratch without them. We ran
that route for every data-matched baseline family on the same deleted rows of the same % SRC: w20-shapes-data-matched/FINDINGS.md
pool, at the baseline's selected schedule, three seeds (Tabs.~\ref{tab:rivals-del}
and~\ref{tab:rivals-del-seeds}, which print each retrain's rows read % SRC: w20-shapes-data-matched/SHAPES-MATCHED.md
and time. After the value and item deletions the n-CLU stays ahead of every % SRC: w20-shapes-data-matched/SHAPES-MATCHED.md
retrained baseline on every seed, with every 95\% interval above % SRC: w20-shapes-data-matched/SHAPES-MATCHED.md
zero. After the red-sphere deletion it is ahead on the mean at every seed, but its interval crosses
zero against the MLP at seeds 1 and 2 and against the LSTM at seed 2 (Fig.~\ref{fig:deletion}): % SRC: w20-shapes-data-matched/SHAPES-MATCHED.md
its own loss on that deletion, about six tuples per held-out combination, is larger than the % SRC: w20-shapes-data-matched/SHAPES-MATCHED.md
change any retrained baseline shows except the continuous MLP at seed 2, so on % SRC: w20-shapes-data-matched/SHAPES-MATCHED.md
collateral damage the baselines' crude retrain is cleaner than the n-CLU's deletion. Nor is the
n-CLU's route faster. Its downdate takes seconds, but the book's refit takes longer % SRC: w20-shapes-data-matched/FINDINGS.md
than any baseline's retrain. What the n-CLU's route has that the retrain does not % SRC: w20-shapes-data-matched/FINDINGS.md; w20-shapes-data-matched/SHAPES-MATCHED.md
% PW2-5 moved -> Tab. rivals-del (rows read, seconds), Tab. rivals-del-seeds (margins, intervals, 26.844 -> 17.275), Tab. del (book refit seconds), Tab. t7 caption (131,072): 131,072; 120,028 to 129,761; 547.7 to 2,064.5 s; +6.954 to +32.607; +2.486; +2.596; +2.477; 26.844 to 17.275; 2,367 to 5,913 s; SRC: w20-shapes-data-matched/SHAPES-MATCHED.md; w20-shapes-data-matched/FINDINGS.md
is the match with a refit without the rows, by construction, and no pool row read. The baselines
retrain at their selected schedule, so the MLP (seeds 0 and 1) and continuous-MLP convergence flags
of Tab.~\ref{tab:t7} apply to them.
% w22 draft-rulings (Ruling 2): the earlier deletion comparators (16,384-row rivals, their retrain-without and the
% NegGrad+ / fine-tune approximate unlearners) are retired from the paper; see w12-deletion-baselines/FINDINGS.md and
% w15-single-design-seeds/T7-3SEEDS.md.

\paragraph{Damaged renders converging to sphere:} A damaged render could be judged sphere for rendering
reasons and score spuriously on the deleted conjunction. On the trained n-CLU's presented here, a degraded render
is judged cube, the held-out shape, so the judged-sphere excess is negative for both maximally
degraded arms ($-0.3291$ for landscape removed, $-0.3078$ for
random atoms) against $+0.0074$ after the red-sphere deletion.

\begin{table}[h]
\centering\small
\caption{Deletion cost for the data-matched baselines: retrain without the rows, on the n-CLU's % SRC: w20-shapes-data-matched/SHAPES-MATCHED.md
131,072-row pool, at the selected schedule, seeds 0 / 1 / 2.}
\label{tab:rivals-del}
\setlength{\tabcolsep}{3pt}
\begin{adjustbox}{max width=\linewidth}\begin{tabular}{lcccccc}
\toprule
baseline, retrain without & seeds & value (red), s & red sphere, s & item (1\%), s & rows read & bytes rewritten \\
\midrule
% SRCBLOCK: w20-shapes-data-matched/SHAPES-MATCHED.md
MLP & 3 & 547.7 / 677.4 / 729.8 & 814.3 / 805.6 / 746.6 & 884.6 / 922.8 / 742.5 & 120,028 / 125,561 / 129,761 & 522,056 \\
LSTM & 3 & 757.5 / 712.4 / 1,016.0 & 772.8 / 767.5 / 1,198.2 & 773.2 / 1,093.8 / 1,034.6 & 120,028 / 125,561 / 129,761 & 522,884 \\
transformer & 3 & 1,298.0 / 1,772.2 / 1,705.3 & 1,255.8 / 1,987.7 / 2,064.5 & 1,246.0 / 1,754.9 / 1,562.8 & 120,028 / 125,561 / 129,761 & 521,732 \\
continuous MLP & 3 & 695.9 / 760.1 / 673.3 & 922.3 / 859.6 / 712.3 & 817.2 / 862.4 / 629.7 & 120,028 / 125,561 / 129,761 & 527,564 \\
% ENDSRC
\bottomrule
\end{tabular}\end{adjustbox}
\end{table}

\begin{table}[h]
\centering\scriptsize
\caption{The n-CLU after deletion against each baseline retrained without the same rows, all % SRC: w20-shapes-data-matched/SHAPES-MATCHED.md
109 combinations, seeds 0 / 1 / 2. Baseline: as trained $\to$ retrained without. n-CLU $-$ baseline: the % SRC: w20-shapes-data-matched/SHAPES-MATCHED.md
paired difference after deletion, with the 95\% bootstrap interval where it crosses zero or comes
closest to it. The n-CLU after deletion reads 50.865 / 50.854 / 52.000 (value, survivor combinations), % SRC: w20-shapes-data-matched/SHAPES-MATCHED.md
41.367 / 40.972 / 42.817 (red sphere) and 47.477 / 47.294 / 49.110 (item).}
\label{tab:rivals-del-seeds}
\setlength{\tabcolsep}{2.4pt}
\begin{adjustbox}{max width=\linewidth}\begin{tabular}{llccc}
\toprule
scope & baseline & seeds & baseline, as trained $\to$ retrained without (0 / 1 / 2) & n-CLU $-$ baseline after deletion (0 / 1 / 2) \\
\midrule
% SRCBLOCK: w20-shapes-data-matched/SHAPES-MATCHED.md
value & MLP & 3 & 39.607 $\to$ 39.124 / 40.449 $\to$ 39.640 / 46.315 $\to$ 44.528 & +11.742 / +11.213 / +7.472 \\
 & LSTM & 3 & 36.685 $\to$ 37.191 / 34.831 $\to$ 35.315 / 43.303 $\to$ 42.742 & +13.674 / +15.539 / +9.258 \\
 & transformer & 3 & 27.933 $\to$ 27.798 / 30.090 $\to$ 29.775 / 40.865 $\to$ 40.157 & +23.067 / +21.079 / +11.843 \\
 & continuous MLP & 3 & 21.393 $\to$ 19.933 / 20.213 $\to$ 26.157 / 29.652 $\to$ 19.393 & +30.933 / +24.697 / +32.607 \\
red sphere & MLP & 3 & 36.706 $\to$ 36.110 / 37.569 $\to$ 38.486 / 43.174 $\to$ 40.220 & +5.257 / \textbf{+2.486 [$-0.94$, 6.02]} / \textbf{+2.596 [$-0.50$, 5.67]} \\
 & LSTM & 3 & 33.972 $\to$ 33.064 / 32.358 $\to$ 33.734 / 40.697 $\to$ 40.339 & +8.303 / +7.239 / \textbf{+2.477 [$-0.87$, 5.82]} \\
 & transformer & 3 & 25.560 $\to$ 24.725 / 27.789 $\to$ 27.349 / 38.651 $\to$ 37.890 & +16.642 / +13.624 / +4.927 [1.29, 8.57] \\
 & continuous MLP & 3 & 17.917 $\to$ 17.697 / 16.541 $\to$ 17.303 / 26.844 $\to$ 17.275 & +23.670 / +23.670 / +25.541 \\
item & MLP & 3 & 36.706 $\to$ 35.881 / 37.569 $\to$ 37.505 / 43.174 $\to$ 42.156 & +11.596 / +9.789 / +6.954 [5.34, 8.66] \\
 & LSTM & 3 & 33.972 $\to$ 33.248 / 32.358 $\to$ 32.587 / 40.697 $\to$ 40.881 & +14.229 / +14.706 / +8.229 \\
 & transformer & 3 & 25.560 $\to$ 25.055 / 27.789 $\to$ 28.083 / 38.651 $\to$ 37.771 & +22.422 / +19.211 / +11.339 \\
 & continuous MLP & 3 & 17.917 $\to$ 19.138 / 16.541 $\to$ 18.450 / 26.844 $\to$ 28.514 & +28.339 / +28.844 / +20.596 \\
% ENDSRC
\bottomrule
\end{tabular}\end{adjustbox}
\end{table}


\subsection{Refusal through the hole:}
\label{sec:hole}

Refusal is a second mechanism on the same object: a coupling table with \emph{holes}.
Punching a hole for a deleted value or conjunction marks the request slot as refused. On the
$\Kstar$ object two holes cost a byte each, read no data rows, take no measurable time and refit nothing, in
a coupling table whose layout is asserted equal to the field-mode n-CLU's over every pair
slot. One hole (red $\wedge$ cube) is refusable although no red cube was ever written: refusal does
not need a deleted row to have existed. % SRC: w13-seeds-and-tables/TABLES.md
The hole is the only mechanism on this object that refuses, and it is a flag: a request for a
holed value is refused by construction, so refusal is measurement (c) Symbolic refusal is not claimed. An item deletion makes nothing impossible: on the n-CLU's own pool, deleting a random
1\% of items clears no value bit and punches no hole, so no request becomes certifiably impossible
and the item row has no refusal row at all. % SRC: w20-shapes-data-matched/SHAPES-MATCHED.json

% ---------------------------------------------------------------------------------------
\section{Light Encoder Adaptation}
\figslot{fig_encoder_adapt}{Encoder adaptation margins: The n-CLU minus the 180k-parameter MLP baseline, in novel tuples per held-out combination, per seed at $\Kstar$, with the frozen encoder,
an adapter and an anchor.}{fig:adapt}

% ---------------------------------------------------------------------------------------
% ---------------------------------------------------------------------------------------
\section{n-CLU is a special case of a Physics motivated CLU}
\label{app:clue}

\paragraph{The derivation:} The context is an orthonormal basis $U_k$ of the
per-well atom span, and the read renders the request's address through that basis,
$x_k=U_k c_k$. The n-CLU's energy for a basis is the distance to its span,
$E_{\text{context},k}=\tfrac12\operatorname{dist}^2(x_k,\operatorname{span}G_k)$, zero everywhere
on the span; the request enters as a basin at the address, $E_{\text{req},k}=\tfrac12\beta\lVert
x_k-a_k\rVert^2$, zero at the address. Because $a_k\in\operatorname{span}(G_k)$ by construction
of the renderer, both terms attain their minima at the same point, so an in-well settle has
nowhere to go. 

\paragraph{The four conditions:} (C1) The n-CLU's energy is distance to the atom span, the atoms
being a basis: with an atoms-as-points energy on the same context the settle moves (relative
displacement median 1.7244, max 2.6682 on a synthetic check). (C2) The address lies in the span:
pushed off the span by $\delta=0.10$ the settle moves by exactly $\delta/(1+\beta)$, measured
0.050000 against the closed form 0.050000. (C3) The state is carried in span coordinates. The same
energy in channel coordinates drifts by \sci{7.447}{-08} relative, float32 rounding, so the paper
says ``identical'' of the channel form and ``bit-identical'' only of the span form. (C4) No
cross-well term: at $\lambda=0$ the displacement is exactly 0.0; at
$\lambda=10^{-4}$, $10^{-2}$, $0.05$, $0.2$ it is \sci{7.875}{-05}, \sci{7.882}{-03},
\sci{3.971}{-02}, \sci{1.720}{-01}. Two wells of the deployed n-CLU were full rank, where an in-well
settle is trivially a no-op. The coercivity certificate reproduces from first principles:
$\lambda^\star=0.25000000$ and a certificate minimum eigenvalue of $-1.50000000$ at the declared
$\lambda=1$, so the landscape is unbounded below at the declared strength and the used strength is
0.225000. % SRC: w12-clue-theorist/CLUE.md
The identity is in-well only; a $\lambda=0$, steps$>0$ arm was never run on the deployed n-CLU (the
synthetic check covers steps 0 / 1 / 10 / 1000); and the \sci{1.9}{-7} off-span residual is the
renderer's float32 rounding, the same size for the static n-CLU that never moves.
% SRC: w12-clue-theorist/CLUE.md


% ---------------------------------------------------------------------------------------

% ---------------------------------------------------------------------------------------
\paragraph{Memory-level deletion in a designed memory:} With an order-independent waitlist, memory-level deletion leaves a memory
byte-equal to one never written with the deleted items, at every load measured, 0.29 to 1.71 times lattice capacity, a membership signal of AUC 0.5000 $\pm$ 0.0000 on all
six statistics and byte equality 1.0000, under three stated conditions (capacity at least the number
of cells, zero leak, priority- or attribute-based eviction. Recency-based eviction is excluded as
intrinsically history-dependent). % SRC: shorts-evidence-inventory.md
The placement rule is the~\citet{blelloch_strongly_2007} stable-matching table with a global priority order;
the theorems follow from theirs and the fix-up cascade is theirs. Its scoping is the one this
paper adopts in Sec.~\ref{sec:object}: an n-CLU-level guarantee only, the frozen encoder and any
residue of past writes in a learned landscape being separate channels, and no exact unlearning can be argued for.

\paragraph{The damping optimum:} On the exact one-step Jacobian the emergent retention
half-life is a V-curve in the damping, with its argmin at $0.902\pm0.003$ of
$\gamma_{\text{crit}}$ and slopes $-1.0020$ and $+1.116$ on three of three emergent seeds, one
damping-optimum curve at two values of the curvature rather than two laws.
% SRC: shorts-evidence-inventory.md
It holds only with fluctuation--dissipation-consistent noise and a Newtonian kinetic term like in the CHLU paper mentioned in Sec.\ref{sec:intro}. On
this study's field-mode n-CLU the rollout misses the same argmin on all three seeds, with the
failing step named: the pure-exponential envelope is invalid at the defective point.
% SRC: w10-physics-section-evidence/EVIDENCE.md


% ---------------------------------------------------------------------------------------
\section{Maze world-model rows}
\label{app:mazewm}


\figslot{fig_cmazet_per_arm}{Each CMAZE-T model succeeding and failing (words plus
observation, test, seed 0). For every model, one held-out $n$-way transition on which its $s_{t+1}$
is correct and one on which it is not, each drawn uniformly with a fixed seed among the qualifying
transitions. The landscape-removed control and the blind-color arm are correct on no held-out
$n$-way transition at seed 0 and show only a failure. The n-CLU has no door success at seeds 0--2
(Fig.~\ref{fig:cmazet-families}), so its success row is the uniform draw. Columns: $s_t$ with the
action, the engine's $s_{t+1}$, and the model's $s_{t+1}$ with the engine's judge's mark. The strip
under each $s_{t+1}$ is read as described in Sec.~\ref{sec:maze-wm}.}{fig:cmazet-per-arm}

\figslot{fig_cmazet_rel_per_arm}{Each CMAZE-T model succeeding and failing with relational
coding (words plus observation, test). As Fig.~\ref{fig:cmazet-per-arm}, for the n-CLU-REL, the
matched baselines, the controls and the blind arms of the same runs, and for MLP-REL, LSTM-REL and
Transformer-REL at 8,000 and 80,000 training steps (the baselines given the same relational columns).
All rows are from seed 0 except the n-CLU-REL success row, which is a door success drawn by rule
among the 44 at seed 2 (Fig.~\ref{fig:cmazet-families}). Every panel is re-read from the runs' % SRC: ../../outputs/pw1-plot-maze/build/a/DOOR-SEARCH.json
persisted models; the re-read reproduces all 132 test scores of seeds 0--5. The strip under each % SRC: ../../outputs/w17-paper/figures/data/cmazet_items_rel.REPRO.json; ../../outputs/pw1-plot-maze/FINDINGS.md
$s_{t+1}$ is read as described in Sec.~\ref{sec:maze-wm}.}{fig:cmazet-rel-per-arm}

\figslot{fig_cmazet_door_then_tint}{Door-then-tint transitions: the move opens a door with
the key in the pocket and enters a tint of the same color (words plus observation, test). Columns as
in Fig.~\ref{fig:cmazet-families}. Rows 1--3: the first such held-out transition at seed 0 for each
move direction present (north, south, west); the n-CLU is wrong on all three, as on every one of its
door-then-tint transitions at seeds 0--2. Rows 4--6: door successes drawn by rule from the n-CLU-REL
at seed 2, one per move direction among its successes (east, south, west; it is wrong on every
northward one), with every column of these rows from that run; fixed-seed draw (seed 16). The strip under each $s_{t+1}$ is read as described in Sec.~\ref{sec:maze-wm}.}{fig:cmazet-door-tint}




\end{document}
